% !PRINTED 338
\weberpage{338}
%<<CONT
in which $f(x)$ and $f''(x)$ have \gesperrt{the same sign}. The
two Figs. 23 and 24 illustrate the situation for positive
$f'(x)$ and positive and negative $f''(x)$.

Thus in the first case $\beta'$, and in the second $\alpha'$, is a better
approximate value than $\beta$ and $\alpha$.

If one wishes at the same time to shift the other limit, so as
to obtain a new narrower interval, one may, following
Fourier, draw from the other end point — that is, in the case of Fig. 23,
at $a$ — the straight line parallel to $b\beta'$, and one obtains the
point $\alpha'$ as the lower limit point of the new interval. Then
in the first case
\[
\alpha' = \alpha - \frac{f(\alpha)}{f'(\beta)}, \quad \beta' = \beta - \frac{f(\beta)}{f'(\beta)},
\]
in the second case
\[
\alpha' = \alpha - \frac{f(\alpha)}{f'(\alpha)}, \quad \beta' = \beta - \frac{f(\beta)}{f'(\alpha)},
\]
and $(\alpha', \beta')$ is a narrower interval in which the sought root lies.

% SOURCE FIGURE CROP: Bd. I registered scan 364 (printed p. 338), 600 dpi,
% crop x=250,y=2250,w=1200,h=1150; source-faithful raster, Fig. 25.
\weberfigureimage{weber-b1-fig25.png}{0.42}

It is hardly necessary to consider
separately the two other cases, in
which $f'(x)$ is negative in the
interval.

But one may also, in order to obtain
two new limits, combine the Newton
method with the interpolation method,
with still better success, as Fig. 25 shows.

One then obtains as the two
new limits:
\[
\begin{gathered}
\alpha' = \alpha - \frac{f(\alpha)(\beta - \alpha)}{f(\beta) - f(\alpha)} = \frac{\alpha f(\beta) - \beta f(\alpha)}{f(\beta) - f(\alpha)}, \\
\beta' = \beta - \frac{f(\beta)}{f'(\beta)}.
\end{gathered}
\]

We summarise the results in the following theorem:

\gesperrt{If an interval $(\alpha, \beta)$ has been marked off, in which
$f(x)$ changes its sign once, and $f'(x)$ and $f''(x)$ not
at all, and if $\beta$ is the end point of the interval
in which $f(\beta)$ and $f''(\beta)$ have the same sign,
no matter whether $\beta$ is smaller or greater than}
%CONT>>
